\documentclass{ximera}

\title{De elektroscoop}
\author{Daan Lipkens}

\begin{document}
\begin{abstract}
    De bouw en werking van een elektroscoop om elektrische ladingen te detecteren.
\end{abstract}
\maketitle

\section{Wat is een elektroscoop?}
Een \textbf{elektroscoop} is een eenvoudig meetinstrument waarmee je kan aantonen of een voorwerp elektrisch geladen is. Daarnaast geeft het een indicatie van de ladingshoeveelheid: hoe meer lading er aanwezig is, hoe sterker de reactie van het toestel.

%---------------------------
% AFBEELDING: Bouw elektroscoop
\begin{figure}[ht!]
    \captionsetup{width=0.85\textwidth}
    \begin{center}
        \includegraphics[width=0.45\textwidth]{pic5/elektriciteit/coulomb/elektroscoop_bouw.jpg}
    \end{center}
    \caption{De bouw van een klassieke bladgoudelektroscoop: een metalen knop, een geleidende staaf en twee dunne metalen blaadjes in een glazen behuizing.}
    \label{fig:elektroscoop_bouw}
\end{figure}
%---------------------------

\section{Bouw en werking}
De klassieke elektroscoop (bladgoudelektroscoop) bestaat uit een glazen of plastic behuizing die de binnenkant beschermt tegen luchtstromingen. Bovenaan zit een metalen knop of schijf die via een geleidende staaf verbonden is met twee zeer dunne metalen blaadjes (vaak goud- of aluminiumfolie) onderaan in de behuizing. De staaf is uiteraard goed geïsoleerd van de behuizing zelf.

\subsection{Werking via contact}
Wanneer je met een geladen voorwerp (bijvoorbeeld een negatief geladen pvc-buis) de metalen knop aanraakt, springen er elektronen over naar de knop. Omdat de knop, de staaf en de blaadjes geleiders zijn, verspreiden deze extra elektronen zich over het hele metalen systeem. 

Hierdoor krijgen beide blaadjes onderaan een \textbf{negatieve} lading. Omdat gelijksoortige ladingen elkaar afstoten volgens de wet van Coulomb, zullen de blaadjes uit elkaar spreiden. Hoe groter de overgedragen lading, hoe verder ze uit elkaar gaan staan.

\subsection{Werking via nadering (polarisatie)}
Zelfs als je de knop \emph{niet} aanraakt, maar het geladen voorwerp er slechts dichtbij houdt, reageert de elektroscoop dankzij influentie. 

Houd je bijvoorbeeld een positief geladen glazen staaf in de buurt van de knop, dan worden de vrije elektronen uit de metalen staaf en de blaadjes naar boven (naar de knop toe) getrokken. Hierdoor ontstaat er onderaan in de blaadjes een tekort aan elektronen. Beide blaadjes worden lokaal \textbf{positief} geladen. Opnieuw stoten de gelijksoortig geladen blaadjes elkaar af en zullen ze spreiden, ondanks dat er geen lading aan het systeem is toegevoegd. Verwijder je de glazen staaf weer? Dan vloeien de elektronen terug en vallen de blaadjes weer netjes naar beneden.

\end{document}